\documentclass[a4paper]{arrowhead}

\usepackage[yyyymmdd]{datetime}
\usepackage{etoolbox}
\usepackage[utf8]{inputenc}
\usepackage[T1]{fontenc}
\usepackage{multirow}
\usepackage{longtable}
\usepackage{tikz}
\usepackage{adjustbox}
\AtBeginEnvironment{longtable}{\small}
\BeforeBeginEnvironment{tikzpicture}{\begin{adjustbox}{max width=\textwidth}}
\AfterEndEnvironment{tikzpicture}{\end{adjustbox}}
\usetikzlibrary{positioning,shapes.multipart,arrows.meta}

\renewcommand{\dateseparator}{-}
\setlength{\emergencystretch}{3em}
\lstset{language=HTTP,breakatwhitespace=false,columns=fullflexible}

\newcommand{\fparam}[1]{\textit{\textcolor{ArrowheadBlue}{#1}}}

%% Special references
\newcommand{\mref}[2][]{{\textcolor{ArrowheadPurple}{\hyperref[sec:model:#2]{\ifstrempty{#1}{#2}{#1}}}}}
\newcommand{\oref}[1]{{\textcolor{ArrowheadBlue}{\hyperref[sec:operations:#1]{#1}}}}
\newcommand{\pdef}[1]{{\textcolor{ArrowheadGrey}{#1 \label{sec:model:primitives:#1} \label{sec:model:primitives:#1s}}}}
\newcommand{\pref}[1]{{\textcolor{ArrowheadGrey}{\hyperref[sec:model:primitives:#1]{#1}}}}

% as template_idd.tex; the input and output arguments are passed already formatted (\mref or \pref)
\newrobustcmd\osubsection[5]{
  \addtocounter{subsection}{1}
  \addcontentsline{toc}{subsection}{\protect\numberline{\thesubsection}operation \textcolor{ArrowheadBlue}{#1}}
  \subsection*{
    \thesubsection\quad
    #2 \textcolor{ArrowheadPurple}{#3} \\
    \small
    \hspace*{0.075\textwidth}\begin{minipage}{0.1\textwidth}
      \vspace*{1mm}
      Operation: \\
      \notblank{#4}{Input: \\}{}
      \notblank{#5}{Output: \\}{}
    \end{minipage}%
    \begin{minipage}{0.825\textwidth}
      \vspace*{1mm}
      \textcolor{ArrowheadBlue}{#1} \\
      \notblank{#4}{#4 \\}{}
      \notblank{#5}{#5 \\}{}
    \end{minipage}
  }
  \label{sec:operations:#1}
}
\newrobustcmd\msubsection[2]{
  \addtocounter{subsection}{1}
  \addcontentsline{toc}{subsection}{\protect\numberline{\thesubsection}#1 \textcolor{ArrowheadPurple}{#2}}
  \subsection*{\thesubsection\quad#1 \textcolor{ArrowheadPurple}{#2}}
  \label{sec:model:#2}
}
%%

\usepackage{xurl}
\begin{document}
\sloppy

%% Arrowhead Document Properties
\ArrowheadTitle{identityManagement HTTP/TLS/JSON}
\ArrowheadServiceID{identityManagement}
\ArrowheadType{Interface Design Description}
\ArrowheadTypeShort{IDD}
\ArrowheadVersion{5.3.0}
\ArrowheadDate{2026-11-09}
\ArrowheadAuthor{Jerker Delsing}
\ArrowheadStatus{REVIEW DRAFT}
\ArrowheadContact{jerker.delsing@ltu.se}
\ArrowheadFooter{\href{www.arrowhead.eu}{www.arrowhead.eu}}
\ArrowheadSetup
%%

%% Front Page
\begin{center}
  \vspace*{1cm}
  \huge{\arrowtitle}

  \vspace*{0.2cm}
  \LARGE{\arrowtype}
  \vspace*{1cm}

  \Large{Service ID: \textit{"\arrowid"}}

  \vspace*{\fill}

  % Front Page Image

  \vspace*{1cm}
  \vspace*{\fill}

  % Front Page Abstract
  \begin{abstract}
    This document describes the HTTP/1.1 protocol with TLS payload security and JSON payload encoding variant of the \texttt{identityManagement} service of the Authentication system.
  \end{abstract}

  \vspace*{1cm}
\end{center}

\newpage
%%

%% Table of Contents
\tableofcontents
\newpage
%%

\section{Documentation edition and implementation scope}
This is the 5.3.0 documentation edition, dated 2026-11-09. It carries forward the system and service contracts from the retained 5.2.2 documentation and identifies the current primary implementation sources. The former \path|documentation/v5.2.2| folder is retained unchanged.

Document version 5.3.0 is distinct from executable, dependency, registered system/service and wire-contract versions. Values still defined by the source (including 5.2.2 artefact versions, 1.0.0 service versions and \path|SysMonitor-AMQP-IDD-5.2.2| payload identifiers) are preserved. This documentation update does not change runtime behaviour or establish a new implementation release. Historical verification statements describe the tests recorded in the inherited baseline; this revision validates document compilation and references.

\section{Overview}
\label{sec:overview}

This document describes the \texttt{identityManagement} service interface of the Authentication system, which provides operator management of the identities known to Authentication (systems, their credentials and sysop flag) and of their active sessions.

It's implemented using protocol, encoding as stated in below table:
\begin{table}[ht!]
  \centering
  \begin{tabular}{|l|l|l|}
    \hline
    \rowcolor{gray!33} Profile type & Type & Version \\ \hline
    Transfer protocol & HTTP & 1.1 \cite{rfc9112} \\ \hline
    Data encryption & TLS & 1.3 \cite{rfc8446}, 1.2 \cite{rfc5246} \\ \hline
    Encoding & JSON & RFC 8259 \cite{rfc8259} \\ \hline
    Compression & N/A & - \\ \hline
    Semantics & N/A & - \\ \hline
    Ontology & N/A & - \\ \hline
  \end{tabular}
  \caption{Communication and semantics details used for the \texttt{identityManagement} service interface}
  \label{tab:comunication_semantics_profile}
\end{table}

This document provides the Interface Design Description IDD to the \textit{identityManagement -- Service Description} document \cite{auth-identityManagement-sd}.
For further details about how this service is meant to be used, please consult that document.

The rest of this document describes how to realize the \texttt{identityManagement} service
interface in detail. Both in terms of its operations (Section \ref{sec:operations}) and its information model (Section \ref{sec:model}).

\subsection{Interface registration}

The service is registered in the ServiceRegistry with the interface template \texttt{generic\_\allowbreak{}https} (\texttt{generic\_\allowbreak{}http} when TLS is disabled), access port 8444 and base path \texttt{/\allowbreak{}authentication/\allowbreak{}mgmt}. The \texttt{operations} property of the interface lists the operations below with their method and path relative to the base path (path parameters are placeholders such as \texttt{\{id\}} that the consumer fills in):

\noindent\begin{longtable}{| p{5.2cm} | p{1.6cm} | p{8.4cm} |} \hline
\rowcolor{gray!33} Operation & Method & Path \\ \hline \endhead
\texttt{identity-\allowbreak{}mgmt-\allowbreak{}create} & POST & \texttt{/\allowbreak{}identities} \\ \hline
\texttt{identity-\allowbreak{}mgmt-\allowbreak{}update} & PUT & \texttt{/\allowbreak{}identities} \\ \hline
\texttt{identity-\allowbreak{}mgmt-\allowbreak{}remove} & DELETE & \texttt{/\allowbreak{}identities} \\ \hline
\texttt{identity-\allowbreak{}mgmt-\allowbreak{}query} & POST & \texttt{/\allowbreak{}identities/\allowbreak{}query} \\ \hline
\texttt{identity-\allowbreak{}mgmt-\allowbreak{}session-\allowbreak{}close} & DELETE & \texttt{/\allowbreak{}sessions} \\ \hline
\texttt{identity-\allowbreak{}mgmt-\allowbreak{}session-\allowbreak{}query} & POST & \texttt{/\allowbreak{}sessions} \\ \hline
\end{longtable}

\newpage

\section{Service Operations}
\label{sec:operations}

The interfaces of the \texttt{identityManagement} service, its operations, data models and implementation are provided below.
In particular, each subsection first names the HTTP method and path used to call the operation, after which it names the abstract operation from the \texttt{identityManagement} SD document, as well as input and output types.

All operations require the requester's identity, presented according to the provider's authentication policy: the client certificate (mutual TLS) and, with the \texttt{outsourced} policy, \texttt{Authorization: Bearer IDENTITY-TOKEN//<token>}; systems holding the Authentication secret key send \texttt{Bearer AUTHENTICATOR-KEY//<system>//<hash>}, and with the \texttt{declared} policy \texttt{Bearer SYSTEM//<system>} is sent. Request and response bodies are \texttt{application/json}.
Errors are answered with an \mref{ErrorMessage} body and the status codes \texttt{400 Bad Request} (invalid parameter), \texttt{401 Unauthorized} (missing or invalid identity), \texttt{403 Forbidden} (not permitted, e.g. by the management or blacklist filter), \texttt{404 Not Found}, \texttt{408 Request Timeout}, \texttt{423 Locked}, \texttt{500 Internal Server Error} and \texttt{503 Service Unavailable} (a system it depends on failed).

\osubsection{identity-mgmt-create}{POST}{/\allowbreak{}authentication/\allowbreak{}mgmt/\allowbreak{}identities}{\mref[Identity\allowbreak{}List\allowbreak{}Mgmt\allowbreak{}Create\allowbreak{}Request]{IdentityListMgmtCreateRequest}}{\mref{IdentityListMgmtResponse}}

Creates identities with credentials of the given authentication method.

A successful call returns \texttt{200 OK} with the output in the body.

\begin{lstlisting}[label={lst:identity-mgmt-create},caption={An \oref{identity-mgmt-create} invocation.}]
POST /authentication/mgmt/identities HTTP/1.1
Host: authentication:8444
Authorization: Bearer IDENTITY-TOKEN//<token>
Content-Type: application/json

{
  "authenticationMethod": "PASSWORD",
  "identities": [
    {
      "systemName": "TemperatureDisplay",
      "credentials": {
        "password": "<password>"
      },
      "sysop": false
    }
  ]
}
\end{lstlisting}

\begin{lstlisting}[label={lst:identity-mgmt-create-resp},caption={An \oref{identity-mgmt-create} response.}]
{
  "identities": [
    {
      "systemName": "TemperatureSensor",
      "authenticationMethod": "PASSWORD",
      "sysop": false,
      "createdBy": "Sysop",
      "createdAt": "2026-10-22T12:00:00Z",
      "updatedBy": "Sysop",
      "updatedAt": "2026-10-22T12:00:00Z"
    }
  ],
  "count": 1
}
\end{lstlisting}

\osubsection{identity-mgmt-update}{PUT}{/\allowbreak{}authentication/\allowbreak{}mgmt/\allowbreak{}identities}{\mref[Identity\allowbreak{}List\allowbreak{}Mgmt\allowbreak{}Update\allowbreak{}Request]{IdentityListMgmtUpdateRequest}}{\mref{IdentityListMgmtResponse}}

Updates credentials and sysop flag of identities.

A successful call returns \texttt{200 OK} with the output in the body.

\begin{lstlisting}[label={lst:identity-mgmt-update},caption={An \oref{identity-mgmt-update} invocation.}]
PUT /authentication/mgmt/identities HTTP/1.1
Host: authentication:8444
Authorization: Bearer IDENTITY-TOKEN//<token>
Content-Type: application/json

{
  "identities": [
    {
      "systemName": "TemperatureDisplay",
      "credentials": {
        "password": "<password>"
      },
      "sysop": false
    }
  ]
}
\end{lstlisting}

\begin{lstlisting}[label={lst:identity-mgmt-update-resp},caption={An \oref{identity-mgmt-update} response.}]
{
  "identities": [
    {
      "systemName": "TemperatureSensor",
      "authenticationMethod": "PASSWORD",
      "sysop": false,
      "createdBy": "Sysop",
      "createdAt": "2026-10-22T12:00:00Z",
      "updatedBy": "Sysop",
      "updatedAt": "2026-10-22T12:00:00Z"
    }
  ],
  "count": 1
}
\end{lstlisting}

\osubsection{identity-mgmt-remove}{DELETE}{/\allowbreak{}authentication/\allowbreak{}mgmt/\allowbreak{}identities?names=\fparam{\{names\}}}{\pref{Array}$<$\allowbreak{}\pref{String}$>$}{}

Removes identities by system name, closing their sessions.

Query parameters: \fparam{names} (List\textless{}String\textgreater{}, required). List values are repeated or comma-separated.

A successful call returns \texttt{200 OK}.

\begin{lstlisting}[label={lst:identity-mgmt-remove},caption={An \oref{identity-mgmt-remove} invocation.}]
DELETE /authentication/mgmt/identities?names=<names> HTTP/1.1
Host: authentication:8444
Authorization: Bearer IDENTITY-TOKEN//<token>
\end{lstlisting}

\osubsection{identity-mgmt-query}{POST}{/\allowbreak{}authentication/\allowbreak{}mgmt/\allowbreak{}identities/\allowbreak{}query}{\mref{IdentityQueryRequest}}{\mref{IdentityListMgmtResponse}}

Returns a page of identities matching the filter.

A successful call returns \texttt{200 OK} with the output in the body.

\begin{lstlisting}[label={lst:identity-mgmt-query},caption={An \oref{identity-mgmt-query} invocation.}]
POST /authentication/mgmt/identities/query HTTP/1.1
Host: authentication:8444
Authorization: Bearer IDENTITY-TOKEN//<token>
Content-Type: application/json

{
  "pagination": {
    "page": 0,
    "size": 10,
    "direction": "ASC",
    "sortField": "createdAt"
  },
  "namePart": "Temp",
  "isSysop": false,
  "createdBy": "Sysop",
  "creationFrom": "2026-10-22T12:00:00Z",
  "creationTo": "2026-10-22T12:00:00Z",
  "hasSession": true
}
\end{lstlisting}

\begin{lstlisting}[label={lst:identity-mgmt-query-resp},caption={An \oref{identity-mgmt-query} response.}]
{
  "identities": [
    {
      "systemName": "TemperatureSensor",
      "authenticationMethod": "PASSWORD",
      "sysop": false,
      "createdBy": "Sysop",
      "createdAt": "2026-10-22T12:00:00Z",
      "updatedBy": "Sysop",
      "updatedAt": "2026-10-22T12:00:00Z"
    }
  ],
  "count": 1
}
\end{lstlisting}

\osubsection{identity-mgmt-session-close}{DELETE}{/\allowbreak{}authentication/\allowbreak{}mgmt/\allowbreak{}sessions?names=\fparam{\{names\}}}{\pref{Array}$<$\allowbreak{}\pref{String}$>$}{}

Closes the sessions of the named systems.

Query parameters: \fparam{names} (List\textless{}String\textgreater{}, required). List values are repeated or comma-separated.

A successful call returns \texttt{200 OK}.

\begin{lstlisting}[label={lst:identity-mgmt-session-close},caption={An \oref{identity-mgmt-session-close} invocation.}]
DELETE /authentication/mgmt/sessions?names=<names> HTTP/1.1
Host: authentication:8444
Authorization: Bearer IDENTITY-TOKEN//<token>
\end{lstlisting}

\osubsection{identity-mgmt-session-query}{POST}{/\allowbreak{}authentication/\allowbreak{}mgmt/\allowbreak{}sessions}{\mref[Identity\allowbreak{}Session\allowbreak{}Query\allowbreak{}Request]{IdentitySessionQueryRequest}}{\mref[Identity\allowbreak{}Session\allowbreak{}List\allowbreak{}Mgmt\allowbreak{}Response]{IdentitySessionListMgmtResponse}}

Returns a page of active sessions.

A successful call returns \texttt{200 OK} with the output in the body.

\begin{lstlisting}[label={lst:identity-mgmt-session-query},caption={An \oref{identity-mgmt-session-query} invocation.}]
POST /authentication/mgmt/sessions HTTP/1.1
Host: authentication:8444
Authorization: Bearer IDENTITY-TOKEN//<token>
Content-Type: application/json

{
  "pagination": {
    "page": 0,
    "size": 10,
    "direction": "ASC",
    "sortField": "createdAt"
  },
  "namePart": "Temp",
  "loginFrom": "2026-10-22T12:00:00Z",
  "loginTo": "2026-10-22T12:00:00Z"
}
\end{lstlisting}

\begin{lstlisting}[label={lst:identity-mgmt-session-query-resp},caption={An \oref{identity-mgmt-session-query} response.}]
{
  "sessions": [
    {
      "systemName": "TemperatureSensor",
      "loginTime": "2026-10-22T12:00:00Z",
      "expirationTime": "2026-10-22T12:00:00Z"
    }
  ],
  "count": 1
}
\end{lstlisting}

\newpage

\section{Data Models}
\label{sec:model}

Here, all data objects that can be part of the service calls associated with this service are listed in alphabetic order.
Note that each subsection, which describes one type of object, begins with the \textit{struct} keyword, which is meant to denote a JSON \pref{Object} that must contain certain fields, or names, with values conforming to explicitly named types.
Fields that are not needed may be omitted or \texttt{null} in requests; empty fields are omitted in responses.
As a complement to the primary types defined in this section, there is also a list of secondary types in Section \ref{sec:model:primitives}, which are used to represent things like hashes, identifiers and texts.

\msubsection{struct}{ErrorMessage}




\noindent\begin{tabularx}{\textwidth}{| >{\raggedright\arraybackslash}p{3.9cm} | >{\raggedright\arraybackslash\small}p{5.3cm} | X |} \hline
\rowcolor{gray!33} Object Field & Value Type & Description \\ \hline
"error\allowbreak{}Message" & \pref{String} &  \\ \hline
"error\allowbreak{}Code" & \pref{Number} &  \\ \hline
"exception\allowbreak{}Type" & \pref{ExceptionType} &  \\ \hline
"origin" & \pref{String} &  \\ \hline
\end{tabularx}



\msubsection{struct}{IdentityListMgmtCreateRequest}

Identities to create.


\noindent\begin{tabularx}{\textwidth}{| >{\raggedright\arraybackslash}p{3.9cm} | >{\raggedright\arraybackslash\small}p{5.3cm} | X |} \hline
\rowcolor{gray!33} Object Field & Value Type & Description \\ \hline
"authentication\allowbreak{}Method" & \pref{String} & Authentication method of the identities; \texttt{PASSWORD} is the only implemented method. \\ \hline
"identities" & \pref{Array}$<$\allowbreak{}\mref{IdentityMgmtRequest}$>$ & Identities to create, update or return. \\ \hline
\end{tabularx}



\msubsection{struct}{IdentityListMgmtResponse}

List of identities.


\noindent\begin{tabularx}{\textwidth}{| >{\raggedright\arraybackslash}p{3.9cm} | >{\raggedright\arraybackslash\small}p{5.3cm} | X |} \hline
\rowcolor{gray!33} Object Field & Value Type & Description \\ \hline
"identities" & \pref{Array}$<$\allowbreak{}\mref{IdentityMgmtResponse}$>$ & Identities to create, update or return. \\ \hline
"count" & \pref{Number} & Total number of matching records (all pages). \\ \hline
\end{tabularx}



\msubsection{struct}{IdentityListMgmtUpdateRequest}

Identities to update.


\noindent\begin{tabularx}{\textwidth}{| >{\raggedright\arraybackslash}p{3.9cm} | >{\raggedright\arraybackslash\small}p{5.3cm} | X |} \hline
\rowcolor{gray!33} Object Field & Value Type & Description \\ \hline
"identities" & \pref{Array}$<$\allowbreak{}\mref{IdentityMgmtRequest}$>$ & Identities to create, update or return. \\ \hline
\end{tabularx}



\msubsection{struct}{IdentityMgmtRequest}

One identity.


\noindent\begin{tabularx}{\textwidth}{| >{\raggedright\arraybackslash}p{3.9cm} | >{\raggedright\arraybackslash\small}p{5.3cm} | X |} \hline
\rowcolor{gray!33} Object Field & Value Type & Description \\ \hline
"system\allowbreak{}Name" & \pref{String} & Name of the system, PascalCase (e.g. \texttt{TemperatureSensor}). \\ \hline
"credentials" & \pref{Object}$<$\allowbreak{}\pref{String}$>$ & Credentials of the authentication method, for \texttt{PASSWORD} the key \texttt{password}. \\ \hline
"sysop" & \pref{Boolean} & Whether the identity is a system operator (default false). \\ \hline
\end{tabularx}



\msubsection{struct}{IdentityMgmtResponse}

A stored identity (without credentials).


\noindent\begin{tabularx}{\textwidth}{| >{\raggedright\arraybackslash}p{3.9cm} | >{\raggedright\arraybackslash\small}p{5.3cm} | X |} \hline
\rowcolor{gray!33} Object Field & Value Type & Description \\ \hline
"system\allowbreak{}Name" & \pref{String} & Name of the system, PascalCase (e.g. \texttt{TemperatureSensor}). \\ \hline
"authentication\allowbreak{}Method" & \pref{String} & Authentication method of the identities; \texttt{PASSWORD} is the only implemented method. \\ \hline
"sysop" & \pref{Boolean} & Whether the identity is a system operator (bypasses the management filters). \\ \hline
"created\allowbreak{}By" & \pref{String} & Name of the system that created the record. \\ \hline
"created\allowbreak{}At" & \pref{DateTime} & UTC time of creation. \\ \hline
"updated\allowbreak{}By" & \pref{String} & Name of the system that changed the record last. \\ \hline
"updated\allowbreak{}At" & \pref{DateTime} & UTC time of the last change. \\ \hline
\end{tabularx}



\msubsection{struct}{IdentityQueryRequest}

Paged filter for identities.


\noindent\begin{tabularx}{\textwidth}{| >{\raggedright\arraybackslash}p{3.9cm} | >{\raggedright\arraybackslash\small}p{5.3cm} | X |} \hline
\rowcolor{gray!33} Object Field & Value Type & Description \\ \hline
"pagination" & \mref{Page} & Paging and sorting of the result. \\ \hline
"name\allowbreak{}Part" & \pref{String} & Part of the system name to match. \\ \hline
"is\allowbreak{}Sysop" & \pref{Boolean} & Only sysop (true) or non-sysop (false) identities. \\ \hline
"created\allowbreak{}By" & \pref{String} & Name of the system that created the record. \\ \hline
"creation\allowbreak{}From" & \pref{DateTime} & Lower bound of the creation time. \\ \hline
"creation\allowbreak{}To" & \pref{DateTime} & Upper bound of the creation time. \\ \hline
"has\allowbreak{}Session" & \pref{Boolean} & Only identities with (true) or without (false) an active session. \\ \hline
\end{tabularx}



\msubsection{struct}{IdentitySessionListMgmtResponse}

List of active sessions.


\noindent\begin{tabularx}{\textwidth}{| >{\raggedright\arraybackslash}p{3.9cm} | >{\raggedright\arraybackslash\small}p{5.3cm} | X |} \hline
\rowcolor{gray!33} Object Field & Value Type & Description \\ \hline
"sessions" & \pref{Array}$<$\allowbreak{}\mref{IdentitySessionResponse}$>$ & Active sessions. \\ \hline
"count" & \pref{Number} & Total number of matching records (all pages). \\ \hline
\end{tabularx}



\msubsection{struct}{IdentitySessionQueryRequest}

Paged filter for active sessions.


\noindent\begin{tabularx}{\textwidth}{| >{\raggedright\arraybackslash}p{3.9cm} | >{\raggedright\arraybackslash\small}p{5.3cm} | X |} \hline
\rowcolor{gray!33} Object Field & Value Type & Description \\ \hline
"pagination" & \mref{Page} & Paging and sorting of the result. \\ \hline
"name\allowbreak{}Part" & \pref{String} & Part of the system name to match. \\ \hline
"login\allowbreak{}From" & \pref{DateTime} & Lower bound of the login time. \\ \hline
"login\allowbreak{}To" & \pref{DateTime} & Upper bound of the login time. \\ \hline
\end{tabularx}



\msubsection{struct}{IdentitySessionResponse}

An active session.


\noindent\begin{tabularx}{\textwidth}{| >{\raggedright\arraybackslash}p{3.9cm} | >{\raggedright\arraybackslash\small}p{5.3cm} | X |} \hline
\rowcolor{gray!33} Object Field & Value Type & Description \\ \hline
"system\allowbreak{}Name" & \pref{String} & Name of the system, PascalCase (e.g. \texttt{TemperatureSensor}). \\ \hline
"login\allowbreak{}Time" & \pref{DateTime} & UTC time of the login. \\ \hline
"expiration\allowbreak{}Time" & \pref{DateTime} & UTC time when the identity token expires. \\ \hline
\end{tabularx}



\msubsection{struct}{Page}

Paging and sorting; all fields optional.


\noindent\begin{tabularx}{\textwidth}{| >{\raggedright\arraybackslash}p{3.9cm} | >{\raggedright\arraybackslash\small}p{5.3cm} | X |} \hline
\rowcolor{gray!33} Object Field & Value Type & Description \\ \hline
"page" & \pref{Number} & Zero-based page number. \\ \hline
"size" & \pref{Number} & Number of records per page (at most \texttt{max.\allowbreak{}page.\allowbreak{}size}, default 1000). \\ \hline
"direction" & \pref{String} & Sort direction, \texttt{ASC} or \texttt{DESC}. \\ \hline
"sort\allowbreak{}Field" & \pref{String} & Field to sort by. \\ \hline
\end{tabularx}



\subsection{Primitives}
\label{sec:model:primitives}

As all messages are encoded using the JSON format \cite{rfc8259}, the following primitive constructs, part of that standard, become available.
Note that the official standard is defined in terms of parsing rules, while this list only concerns syntactic information.
Furthermore, the \pref{Object} and \pref{Array} types are given optional generic type parameters, which are used in this document to signify when pair values or elements are expected to conform to certain types.

\noindent\begin{tabularx}{\textwidth}{| p{4.6cm} | X |} \hline
\rowcolor{gray!33} Type & Description \\ \hline
\pdef{Array} & A JSON Array \cite{rfc8259}; \pref{Array}$<$A$>$ has elements of type A. \\ \hline
\pdef{Boolean} & JSON Boolean, \texttt{true} or \texttt{false}. \\ \hline
\pdef{DateTime} & JSON String with an ISO 8601 UTC timestamp, e.g. \texttt{2026-10-22T12:00:00Z}. \\ \hline
\pdef{ExceptionType} & One of \texttt{ARROWHEAD}, \texttt{INVALID\_\allowbreak{}PARAMETER}, \texttt{AUTH}, \texttt{FORBIDDEN}, \texttt{DATA\_\allowbreak{}NOT\_\allowbreak{}FOUND}, \texttt{TIMEOUT}, \texttt{LOCKED}, \texttt{INTERNAL\_\allowbreak{}SERVER\_\allowbreak{}ERROR}, \texttt{EXTERNAL\_\allowbreak{}SERVER\_\allowbreak{}ERROR} (JSON String). \\ \hline
\pdef{Number} & JSON Number, an integer. \\ \hline
\pdef{Object} & JSON Object; \pref{Object}$<$A$>$ has string keys and values of type A. \\ \hline
\pdef{String} & JSON String. \\ \hline
\end{tabularx}

\newpage

\section{Go and Rust implementation variants}
\subsection{Go implementation}
Primary system folder: \path|Eclipse-Arrowhead-go/ah5-core-go/authentication|.\par
The system owns its Go module, launcher and build configuration. Use the per-system \path|Makefile| and \path|go.mod|; the containing Go tree is the build source, not a historical Arrowhead copy. Ordinary native roles use independent identities and JSON snapshot state rather than the Java database. Do not assume Java SQL migration, clustering or identical transaction/retry semantics. The registry has its own configuration and optional native AMQP monitor; other roles must be checked against their own registrations.
\paragraph{Listening port.} Built-in default 127.0.0.1:8444 (Arrowhead default). The shipped config.json sets "listen": "0.0.0.0:8444". Override with "listen" in the configuration file, SERVER\_PORT=N or -server.port N; the port overrides keep the configured host.
Native design/route evidence: \path|Eclipse-Arrowhead-go/ah5-core-go/authentication/documentation|.\par
For this Java interface, select only a native transport and operation explicitly implemented and advertised by that port. A matching method/path does not establish DTO, error, authorization, retry or persistence equivalence. Java MQTT/AMQP bindings are not automatically provided by a Go or Rust role. Differential contract and security checks remain necessary before substitution.
\subsection{Rust implementation}
Primary system folder: \path|Eclipse-Arrowhead-rust/ah5-core-rust/authentication|.\par
The per-system entry point is built by the \path|arrowhead-core-rust| crate. Build from the primary Rust workspace with the applicable package/binary selection. Documentation links can resolve into the shared core/support catalogue. Independent process identity and state remain required.\par
The ordinary native profile uses Tokio/Axum and rustls with peer-identity checks and independent locked JSON state. Network effects are not distributed database transactions. The precise trust, CRL reload and persistence guarantees are those of the native profile, not the Java implementation.\par
Native design/route evidence: \path|Eclipse-Arrowhead-rust/ah5-core-rust/documentation/Authentication|.\par
For this Java interface, select only a native transport and operation explicitly implemented and advertised by that port. A matching method/path does not establish DTO, error, authorization, retry or persistence equivalence. Java MQTT/AMQP bindings are not automatically provided by a Go or Rust role. Differential contract and security checks remain necessary before substitution.

\bibliographystyle{IEEEtran}
\bibliography{bibliography}

\newpage

\section{Revision History}
\subsection{Amendments}

\noindent\begin{tabularx}{\textwidth}{| p{1cm} | p{3cm} | p{2cm} | X | p{4cm} |} \hline
\rowcolor{gray!33} No. & Date & Version & Subject of Amendments & Author \\ \hline

1 & 2026-10-22 & 5.2.2 & Document created from the Arrowhead LaTeX templates for the 5.2.2 Java Spring implementation, including the local extensions of this workspace. & Jerker Delsing \\ \hline
2 & 2026-11-09 & 5.3.0 & New documentation edition; primary-source references and Go/Rust implementation scope; 5.2.2 retained. & Documentation maintenance \\ \hline
\end{tabularx}

\subsection{Quality Assurance}

\noindent\begin{tabularx}{\textwidth}{| p{1cm} | p{3cm} | p{2cm} | X |} \hline
\rowcolor{gray!33} No. & Date & Version & Approved by \\ \hline

1 & & \arrowversion  &  \\ \hline

\end{tabularx}

\end{document}
